The \NAME{} auto-scheduler finds high-performance schedules for an input program without
user intervention. \NA{Three challenges shape its design.
  First, advanced fusion is non-local: it fuses two reductions together with all statements
  on the path between them, so it cannot be decided one statement at a time;
  the auto-scheduler looks ahead along data dependencies to find these opportunities
  (\S\ref{sec:adv-fusion}).
  Second, fusion can leave the program in forms that tile compilers optimize poorly,
  such as scan dependencies;
  the auto-scheduler regularizes the program to keep it tile-ready (\S\ref{sec:regularize}).
  Third, the best schedule and tile compilers depend on the operator, input shape, and GPU;
  the auto-scheduler sketches a set of promising schedules,
  which we call \textit{schedule seeds},
  and leaves their decision parameters to the auto-tuner (\S\ref{sec:auto-tuner}),
while heuristics prune the seeds aggressively to keep auto-tuning \mbox{within minutes.}}

\subsection{Auto-scheduling Algorithm}\label{sec:autosched-alg}
The auto-scheduler builds up schedule seeds by applying a sequence of \emph{scheduling passes}
(Table~\ref{tab:passes}), each applying one class of optimizations such as tiling or fusion.
% The input to the auto-scheduler is a program in which all compute statements
% are perfectly nested in their own loops.
It visits compute statements in a pre-order traversal, producers to consumers,
while maintaining a list of candidate schedules.
At each statement, it applies the passes in order to every candidate;
a pass may replace a candidate with several if there are multiple promising ways
to transform the statement, keeping only those its heuristics favor
(e.g., fusing into the deepest feasible loop).
After the last statement, the final two passes run on each finished candidate,
and the candidates become the schedule seeds,
whose design parameters (e.g., block tiling sizes) the auto-tuner explores.

\begin{table*}[t]
  \centering
  \small
  \caption{Compile passes in the Nautilus auto-scheduler, in the order they are applied.
    The last two run on
  each finished candidate.}
  \label{tab:passes}
  \begin{tabular}{lll}
    \hline
    Scheduling Pass & Transformations & Description \\
    & Used & \\
    \hline
    PadRows            & 1                & Adds a candidate with short query rows
    (decode, MTP) padded to 16 for tensor cores \\
    InlineReshapes     & 1                & Inlines tensor reshaping \\
    AdvTilingFusion    & 4   & Uses advanced Neptune fusion to fuse compute statements \\
    CacheSetup         & 4  & Shares or rematerializes large intermediates per tile in
    shared memory \\
    DecomposeSplitScan & 3      & Binds grid loops; apply decompose reduction or split
    scan buffer \\
    & & to every compute statement which had advanced fusion \\
    BackendSelection   & (none)            & Configures compile parameters for different backends \\
    BlockOrder         & 1                & Makes the order of statements in each loop a
    tuning choice \\
    DataPlacement      & 4   & Places every tensor in global, shared memory or registers;
    caches index values \\
    Finalize           & 2          & Turns masks into loop ranges and folds accumulators
    before lowering to tile IR \\
    \hline
  \end{tabular}
\end{table*}

\subsection{Auto-Scheduler Transformations}
Our auto-scheduler applies four families of transformations:
tiling, fusion, cache management, and dataflow management.
These transformations capture the main steps needed to translate
and optimize a tensor-expression operator into a tile program ready for \NAME{}'s tile optimizers.
% A family of transformations may be implemented as multiple passes in the auto-scheduler
% to break dependencies between certain transformations. \matei{It might be unnecessary
% to talk about
% passes at all if we give an idea of the transformation types.}

\subsubsection{Tiling.}
The bi-level tiling transformation applies to a perfect loop nest,
splitting each loop into an inner loop and an outer loop,
and reordering them so that all outer loops are outside of all inner loops.
For example, a loop nest with two loops \texttt{i, j} turns into the loops
\texttt{i0, j0, i1, j1}, with the first two loops being the outer loops.
The other passes in the auto-scheduler transform only the outer loops;
the inner loops remain perfectly nested so they cleanly map to tiles in the tile IRs later.
Later, outer loops that are parallelizable (no data dependency across iterations)
are mapped to GPU blocks.
A size limit on inner loops is also enforced to prevent tiles from becoming
too large and not fitting into the cache - in this case, the respective large inner loop is itself
tiled into an inner and outer loop if possible.

% \yifan{Good to have but I'm struggling where to put this, maybe not here.}
% \matei{moved to dataflow}

\subsubsection{Operator Fusion.}\label{sec:adv-fusion}
Fusion reduces the number of kernels and improves data locality.
Each fusion primitive fuses a compute statement into a loop nest.
\textbf{Classic loop fusion} preserves data dependencies
and fuses a statement into the element-wise loop nest of a direct producer.
Fusing into a reduction loop nest requires \textbf{advanced fusion}~\cite{neptune},
which fuses reductions into an ``online'' form by intentionally breaking some dependencies,
and repairing the program with changed compute expressions, at the cost of extra computation.
Advanced fusion applies to a pair of reductions,
but also needs to capture all element-wise statements on the path between them.

\NAME{} has two advanced fusion primitives.
\textbf{Rolling update} fuses the reductions into a single kernel.
\textbf{Split-k update} computes a partial result for each iteration of the reduction loop
in parallel thread blocks, then combines the partials in a second kernel with the same repair;
this gains parallelism along the reduction axis when the other loops cannot fill the GPU,
as happens in decoding attention.
The auto-scheduler is capable of \textbf{hybrid strategies}, applying split-k update across chunks
of the reduction and rolling update within each chunk.
% which suits kernels with few query rows but a long reduction,
% such as multi-token prediction verification.
% The auto-scheduler chooses the strategy when tiling the loop nest.

For each program statement $s$, the fusion pass first greedily applies classic fusion
into the loop nests of its direct producers, at the deepest possible level.
% \yifan{I moved this over from ``additional considerations'' -- it can be beneficial, but
% what do we actually do?} \matei{not supported yet to fuse with multiple producers, we try
% either one as a candidate}
When a statement has multiple producers, the pass tries fusing with each one and produces
multiple candidate schedules.
If $s$ is a reduction, the pass also tries advanced fusion into each predecessor
reduction, which, when feasible, always fuses deeper than classic fusion.
If an element-wise $s$ cannot be fused, the pass looks ahead for a successor reduction
whose advanced fusion would capture $s$; if none exists, $s$ starts a new kernel.

Consider a loop nest $L$ with an element-wise loop $e$ around a reduction loop $r$.
$L$ has an element-wise consumer $c_1$, and $c_1$ has a reduction consumer $c_2$.
On its own, $c_1$ can only fuse into $e$ via classic fusion.
Advanced fusion of $c_2$ into $r$ would also capture $c_1$ and place it under $r$,
deeper than classic fusion can, but this opportunity is only visible when examining $c_2$;
this is what makes advanced fusion \emph{non-local}.
Since pre-order traversal visits $c_1$ first, the fusion pass greedily fuses $c_1$ into $e$.
This does not prevent the deeper fusion: when the pass later fuses $c_2$ into $r$,
it fuses $c_1$ again, moving it from under $e$ to under $r$.
Therefore, \NAME{}'s fusion pass handles non-local advanced fusion with a plain
pre-order traversal, without anticipating advanced fusion when making classic fusion decisions.

% Matei: moved earlier | and a ``rolling update'' technique that enables fusion between reductions.
% Rolling update is an advanced fusion technique from Neptune \cite{neptune}
% which is beyond the capabilities of classic loop fusion,
% fusing multiple reductions into an ``online'' kernel.
% While classic loop fusion preserves data dependency,
% rolling update intentionally breaks some dependency constraints
% and repairs the program by changing compute expressions.
% \matei{the greedy heuristic part might make more sense to include in section 4.4
% instead of here if we are discussing advanced fusion there}
% and in the case of RollingUpdate, we prefer to capture more loop nests.
% In the case where there isn't a Pareto-dominant choice (such as when one fusion choice
% is deeper but
% the other captures more loop nests), both strategies are considered as produced schedules.
% The fused loop nests do not have to be
% adjacent, in which case any element-wise loop nests along the data dependency path
% between the target loop nests
% will be `captured' and fused as well
% This pass operates in a pre-order traversal of the schedule. The program's compute statements are
% traversed from producer to consumer, and existing candidate schedules are evaluated for
% fusion opportunities
% at the compute statement - when fusion is possible with one of that compute statement's loops,
% new candidate schedules are added, replacing the original unfused candidate.

\subsubsection{Cache Management.}
Cache management places tensors on an inner (more local) memory level,
i.e., shared memory or registers rather than global HBM, to reduce slow global memory traffic.
All tensors start in global memory. The auto-scheduler performs two types of data
localization: (1) it caches a copy of a tensor to shared memory if \NA{a reduce loop
nest} exhibits temporal reuse of the tensor, and (2) it places the tensor itself
in an inner memory level if all of its accesses are local to a thread block.
Tensors that are not explicitly cached are placed by how they are used:
in global memory if another kernel reads them (such as split-k partials),
\NA{in shared memory for matmul operands whose layout is incompatible with
  tensor core matmul, and in registers otherwise.
  A tensor core matmul has strict requirement for the layout of its register operands,
  which can cause tile compilers to fail on programs.
  \NAME{} anticipates layout failures using heuristic rules and places tensors on shared memory
when a relayout is needed.}
% \NA{ A matmul's second
%   operand always goes to shared memory: every warp reads all of its columns, so keeping it
%   in registers would require redistributing it across threads before each matmul. The first
%   operand may be the result of an earlier matmul, held in registers and spread across
%   threads in that matmul's output layout. It can stay there only if this layout matches the
%   operand layout the second matmul expects, which holds under the following conditions: (1)
%   The second matmul reads the tile along its reduction axis. (2) The tile has at least 16
%   rows per warp (128 rows for the 8 warps we can sample), so both matmuls split rows across
%   warps the same way. With fewer rows, threads hold elements the consumer expects in other
%   threads, and moving them would need a cross-thread exchange that the tile optimizer does
%   not insert. (3) Finally, the operand is not FP8. FP8 matrix instructions pack four
%   consecutive reduction-axis elements per register, while the FP32 accumulator gives each
% thread pairs of adjacent columns, so the layouts will never match.} \matei{this may be
% too detailed...}
% Later in the tile optimizer, tensors whose register layout does not match what their
% consumer expects are rejected, so cache management places these tensors in shared memory
% such that the tile optimizer can reload them with a relayout.
% ; this applies to a matmul's second operand, and an 8-bit operand computed in registers,
% which inherits its FP32 accumulator's layout while FP8 tensor cores expect a different one.

Inner memory levels such as shared memory are small,
so the auto-scheduler analyzes tensor sizes to identify
large buffers with multiple readers that would otherwise occupy scarce local memory.
If such a buffer's producer and all its readers share a loop, and the producer's
tile at that loop fits in shared memory, the auto-scheduler keeps one copy of
that tile for all readers. Otherwise, it rematerializes the buffer,
recomputing the buffer near each use to reduce memory pressure and improve occupancy.
Integer tensors used for indexing, such as sparse attention's top-k indices,
are placed by how they are accessed.
\NAME{} tries to eliminate integer intermediates used only for indexing (in a gather
operation) by inlining their computation into their consumers.

% Each thread then loads the indices it needs directly from global memory, instead of being
% restricted to those held in its own registers, so the gather can be spread across all warps.
% \NA{When the auto-scheduler folds accumulators,} an integer input used both in gathers
% and as a value in computation gets a shared memory copy for the value consumer only,
% staged once per tile and prefetched with the gathers, while the gathers still read it
% directly from global memory.

\subsubsection{Dataflow Management.}\label{sec:regularize}
Our auto-scheduler applies several transformations to align the data dependency patterns
of the program with what tile compilers expect,
as tile compilers often optimize only a small set of tile computation patterns well.
One pattern they handle poorly is the ``scan'' dependency
(e.g., \texttt{a[i, j] = f(a[i - 1, j])}),
which can exist in the input program itself
or arise from previous transformations, such as rolling update;
existing tile compilers often fail to optimize scan dependencies
or produce significantly suboptimal kernels.
The auto-scheduler removes them with a split-scan-buffer transformation,
which multi-buffers the scan's output and updates the buffers once per loop iteration.
\NAME{} uses range analysis to convert expressions of boolean masks
into tightened loop ranges that skip fully masked tiles,
which improves the performance of masked attention, such as causal attention.
\NA{Because matmuls with fewer than 16 rows cannot use tensor cores,
  \NAME{} also adds candidates that pad such row axes to 16 (e.g., query rows in decoding),
zero-filling the computation and masking the load and stores.}

% To maintain correctness, if a value computed inside a rolling update loop is read
% from outside of that loop, the candidate is rejected, as the value would not be properly rescaled.
% \yifan{Sentence moved from Auto-Tuner.}

\setminted{breaklines=true, breakanywhere=true, fontsize=\footnotesize, fontfamily=lmtt,
xleftmargin=0pt, xrightmargin=0pt, framesep=1pt}
\begin{figure*}[ht]
  \begin{subfigure}[b]{0.36\textwidth}
    \centering
    \begin{minted}{python}
for bx in range(N0, "blockIdx.x"):
  for i1 in range(N1):
    xmax[bx*N1+i1] = -inf
  for j0 in range(M0):
    for i1, j1, k in grid(N1, M1, K):
      i, j = bx*N1+i1, j0*M1+j1
      mm[i, j] += Q[i, k] * K[j, k]
    for i1, j1 in grid(N1, M1):
      i, j = bx*N1+i1, j0*M1+j1
      xmax[i] = max(
        xmax[i], mm[i, j])
    \end{minted}
    \caption{Scalar IR, with scalar memory accesses and computation under loop nests.
      This IR facilitates high-level optimizations such as operator fusion,
    which are central to \NAME{}'s auto-scheduler.}
  \end{subfigure}
  \hspace{2pt}
  \begin{subfigure}[b]{0.32\textwidth}
    \centering
    \begin{minted}{python}
bx = block("blockIdx.x", 0, N0)
vmax_init = tile(-inf, [N1])
vmax = for j0 in loop(M0)
    params (vmax_1 = vmax_init) {
  il, ir = bx*N1, (bx+1)*N1
  jl, jr = j0*M1, (j0+1)*M1
  vmm = dot(Q[il:ir, 0:K],
    K[jl:jr, 0:K].T)
  vmax_local = max(vmm, axis=1)
  vmax = maximum(
    vmax_local, vmax_1)
  yield vmax
}
store(vmax, xmax[bx*N1:bx*N1+N1])
    \end{minted}
    \caption{Virtual register (VR) tile IR facilitates expression rewrites.
    Tiles are expressed as values instead of tensor accesses.}
    \label{fig:vr-ir-example}
  \end{subfigure}
  \hspace{2pt}
  \begin{subfigure}[b]{0.31\textwidth}
    \centering
    \begin{minted}{python}
bx = block("blockIdx.x", 0, N0)
for j0 in range(M0):
  il, ir = bx*N1, (bx+1)*N1
  jl, jr = j0*M1, (j0+1)*M1
  xmm[0:N1, 0:M1] = dot(
    Q[il:ir, 0:K],
    K[jl:jr, 0:K].T)
  xmax_local = tile(-inf, [N1])
  xmax_local[0:N1] = max(
    xmm[:, :], axis=1)
  xmax[0:N1] = maximum(
    xmax_local[:], xmax[:])
    \end{minted}
    \caption{Memory access (MA) tile IR desugars the VR-tile IR and exposes data
      movement between tensors. \NAME{} uses this IR to interface with its tile
      optimizers (Triton, Tawa, and TileLang).
    }
  \end{subfigure}
  \vspace*{-1.5em}
  \caption{An example program in \NAME{}'s three IRs.}
  \label{fig:3ir-example}
  \vspace*{-1em}
\end{figure*}


\subsection{Auto-Tuner}\label{sec:auto-tuner}
\NAME{} uses MetaSchedule's auto-tuner~\cite{metaschedule} to search the decision parameters
of the schedule seeds.
Every tunable decision is a sampled variable in the schedule trace,
which MetaSchedule's evolutionary search mutates and child schedules inherit:
the search generates schedules by varying these decisions, predicts their performance
with a learned cost model, and measures the projected-best schedules on the target GPU.
The tuning budget is the number of GPU measurements.
% \matei{The search space is pruned by dropping any seeds which fail to build.
%   Many fusion structures can't be built with any tile sizes - these would otherwise dominate
%   the initial tuning population.
%   \textit{note: is this something to include? might make people question why we are
% producing invalid seeds?}}

\NA{\textbf{Tunable decisions.}
  A full list of tunable decisions is given in Appendix Table~\ref{tab:decisions}.
  Structural decisions become separate seeds during auto-scheduling (\S\ref{sec:autosched-alg}),
  such as fusion strategy, whether row padding is applied, and the tile compiler to use.
  The BackendSelection pass splits each seed into one per tile compiler (Triton,
  Tawa, or TileLang).
  Within each seed, there are integer and categorical knobs to be tuned.
  Some apply directly to the schedule, such as tile sizes and the KV split count.
  The order of certain expensive statements (matmuls, reductions, copies) is also tunable.
  Other knobs become annotations on the IR program and take effect in later stages:
  switchable expression rewrites (such as accumulator folding),
  and tile optimizer parameters (e.g., TMA copy flags).
  % BackendSelection records these as annotations that the implementing stage reads.
  This design allows the auto-tuner to explore all these decisions from multiple levels,
jointly making decisions on tiling, fusion, expression rewrite, and tile optimization.}
% \matei{determine if table will be included}}

\section{\NAME{}'s Successive Lowering Pipeline}
\label{sec:opt-pipeline}

After scheduling, \NAME{} lowers the program from its scalar IR through two tile IRs,
because no single IR suits both remaining categories of
optimizations: expression rewriting needs tiles as values that it can freely combine and move
across statements and loop iterations, while tile compilers need explicit tile loads and stores
to perform SIMD optimizations.
The virtual register tile IR (VR-tile IR) serves the former and is new in \NAME{};
the memory access tile IR (MA-tile IR) serves the latter and is purposefully similar to
the tile IRs of existing tile compilers, simplifying interfacing with them.

\textbf{Tile IRs.}
Figure \ref{fig:3ir-example} shows an example program in both tile IRs,
and scalar IR for reference.
The MA-tile IR exposes tile data movement as explicit tensor accesses,
which is typical in tile IRs of existing tile compilers.
\NAME{}'s VR-tile IR expresses tiles as local variables (``virtual registers'')
rather than as memory accesses, similar in spirit to memory-to-register promotion
(``mem2reg'') in compilers such as LLVM.
This design reduces data movement instructions that would otherwise obstruct expression
manipulation, and enable powerful rewrites.
For example, \texttt{vmm}, \texttt{vmax\_local}, and \texttt{vmax} in Figure \ref{fig:vr-ir-example}
are each bound to a tile expression, so expression rewrites view them as one
expression when needed and freely move terms across them.

Virtual registers cover values within a loop body, but a value that evolves across loop
iterations, such as a running maximum, would still be expressed as repeated update-writes
to a tensor, hiding it from expression rewrites.
The VR-tile IR therefore includes a for-loop expression, which takes loop-carried
parameters and returns values: in Figure \ref{fig:vr-ir-example}, the loop expression
\texttt{for j0 in loop(M0)}
takes parameter \texttt{vmax\_1} and returns \texttt{vmax}.
Loops with loop-carried values also exist in other IRs, e.g., MLIR's \texttt{scf.for}~\cite{mlir};
the VR-tile IR uses them over tiles, so that tile value evolution over loop iterations
is expressed natively and tensor expression rewriting can operate across iterations.

\textbf{Lowering to VR-tile IR.}
\NAME{} lowers the scheduled scalar IR into the VR-tile IR,
relying on the tile-ready form that the auto-scheduler maintains (\S\ref{sec:auto-sched}).
Each perfectly nested inner loop nest becomes a tile expression bound to a virtual register,
such as the \texttt{dot} for the matmul in Figure~\ref{fig:vr-ir-example};
an outer loop that updates a tensor across its iterations, such as the running maximum
\texttt{xmax}, becomes a loop expression whose loop-carried parameter (\texttt{vmax\_1})
holds that tensor's tile.

\textbf{Expression Rewriting.}
\NAME{}'s expression rewriter applies a set of rewrite rules recursively on the VR-tile IR.
Table \ref{tab:rewrite-rules} lists representative rules, including rules over loop expressions
and rules informed by GPU hardware (e.g., rewriting $e^x$ using $2^x$).
These rules matter more after advanced fusion,
whose repair computation (\S\ref{sec:adv-fusion}) introduces rescaling in every loop iteration.
For example, the fourth rule in Table \ref{tab:rewrite-rules}
cancels out division and multiplication over loop iterations, removing $2n - 1$
operations from the program,
and the second rule computes the online-softmax rescale once per row instead of per element.
Expression rewriting reduces the latency of \NAME{}'s kernels in our ablation study
(Table~\ref{tab:ablation}).

\begin{table}[ht]
  \centering
  \scriptsize
  \setlength{\tabcolsep}{3pt}
  \renewcommand{\arraystretch}{1.1}
  \caption{Representative rules in \NAME{}'s expression rewriter. Standard algebraic
  identities are used but not listed.}
  \tt
  \label{tab:rewrite-rules}
  \begin{tabular}{l}
    \hline
    \textbf{Pattern} $\Rightarrow$ \textbf{Replacement} \\
    \hline
    $ \exp(x) \Rightarrow \mathrm{exp2}(\mathrm{log2}(e) \cdot x) $ \\[2pt]
    $\mathrm{exp2}(\text{bcast}(a) - \text{bcast}(b))$ $\Rightarrow$
    $\text{bcast}(\mathrm{exp2}(a - b))$ \\[2pt]
    $ x = \text{for}(...)\ \{ \text{yield } c \cdot x \} $ \\
    \quad $\Rightarrow$ $ x' = \text{for}(...)\ \{ \text{yield } x \};\ x = c \cdot x' $ \\[2pt]
    $ x = \text{for}(i: n)\ \{ \text{yield } x / y_i \cdot y_{i+1} \} $ \\
    \quad $\Rightarrow$ $ x' = y_n \cdot \text{for}(i: n)\ \{ \text{yield } x \} $ \\[2pt]
    $\text{select}(c, \text{dot}(A,B) \cdot s, k)$ \\
    \quad $\Rightarrow$ $\text{dot}(A, B, \text{init} = \text{select}(c, 0, k/s))\cdot s$ \\
    \hline
  \end{tabular}
\end{table}

\textbf{Tile Optimizations.}
After expression rewriting, \NAME{} translates the program into the MA-tile IR
by desugaring loop expressions and tile-valued virtual registers
into plain for-loops and explicit tile loads and stores.
The MA-tile IR builds on the tile IR of Neptune, which Neptune lowers to Triton.
We design the MA-tile IR as a superset of multiple existing tile languages,
which enables \NAME{} to use multiple tile compilers as its tile optimizers
(Triton, Tawa~\cite{tawa}, and TileLang~\cite{tilelang}), as chosen by the auto-tuner
(\S\ref{sec:auto-tuner}).
